\section{Projektumfeld und technische Umgebung}

\subsection{Technische Umgebung}

Die Entwicklung erfolft im Unternehmensumfeld und deren Systemen.
GitHub wird zur Bereistellung der Lern-Repositories verwendet und dient der Versionierung der Inhalte.
\\
Die Skripte werden in Bash und Python umgesetzt, auf einer Linuxumgebung entwickelt und ausgeführt.
Für die Schnittstelle zwischen Datenbank und Frontend wird eine API eingesetzt mit dem FastAPI Framework aus Python.
\\
Neo4j dient als Graphdatenbank zur Speicherung der strukturierten Aufgaben(Challenges) und ihren Beziehungen(Abhängigkeiten).
\\
Die API und Neo4j werden in einer Docker-Umgebung in einem Container auf Linux betrieben und über eine Docker-Compose miteinander verbunden.
\\
Das Frontend bekommt die Daten über die API und besitzt keinen direkten Zugiff auf die Neo4j-Datenbank.

\subsection{Systemarchitektur und Datenfluss}

Der Datenfluss des Systems ist wie folgt aufgebaut:

\begin{verbatim}
                                GitHub-Repositories
                                        |
                                Bash-/Python-Skripte
                                        |
                                  graph-data.json
                                        |
                                      Neo4j
                                        |
                                     FastAPI
                                        |
                                     Frontend
\end{verbatim}
Die Projektdokumentation wird dahingegen in LateX erstellt.
\\
Der Build der Dokumentation wir durch eine Makefile automatisiert.
Zudem erledigt ein GitHub-Workflow in einer Tag/Release-Versionierung, den automatisierten Build und die Ablage des erzeugten Dokuments.
